\documentclass[11pt]{article}
\usepackage[utf8]{inputenc}
\usepackage[T1]{fontenc}
\usepackage{booktabs}
\usepackage{amsmath}
\usepackage{geometry}
\geometry{margin=2.5cm}

\title{Causal Fairness Analysis Report}
\author{Generated by FairMind and an LLM}
\date{2026-08-22}

\begin{document}
\maketitle

\section{Analysis setup}

\begin{tabular}{ll}
\toprule
Dataset & adult \\
Protected attribute ($X$) & S2\_gender \\
Comparison & Female $\to$ Male \\
Outcome ($Y$) & T\_income (>50K) \\
Mediator ($W$) & hours-per-week \\
Confounder ($Z$) & education \\
\bottomrule
\end{tabular}

\section{Causal effects (FairMind, exact values)}

The five effects below are computed by FairMind through exact inference on the
fitted Bayesian network. These values are final and must not be recomputed or
altered.

\begin{tabular}{lr}
\toprule
Effect & Value \\
\midrule
Total Variation (TV) & 0.193700 \\
Total Effect (TE) & 0.184852 \\
Spurious Effect (SE) & 0.008848 \\
Direct Effect (DE) & 0.184000 \\
Indirect Effect (IE, decomposition form: TE = DE + IE) & 0.000852 \\
\bottomrule
\end{tabular}

\section{Qualitative interpretation}

\subsection{Total effect and spurious component (TV, TE, SE)}

The total disparity between the two groups is significant, with a total variation of 0.1937. After removing the confounding effect of education, the genuine causal effect (TE) remains substantial at 0.184852. The spurious effect driven by the confounder education is relatively small at 0.008848, indicating that the majority of the observed disparity is due to the direct and indirect causal effects rather than confounding.

\subsection{Direct effect (DE)}

The direct effect of gender on income is 0.1840, indicating a strong direct discrimination against the protected group (females). This direct effect is very close to the total effect, suggesting that the direct pathway from gender to income is the primary driver of the disparity.

\subsection{Indirect effect (IE)}

The indirect effect through the mediator hours-per-week is 0.000852, which is negligible. Since the indirect effect has the same sign as the direct effect, it slightly amplifies the direct effect, but the overall impact is minimal.

\section{Recap Questions}

\begin{enumerate}
\item Is there direct discrimination against the protected group? \textbf{Answer:} YES
\item Does the indirect effect through the mediator mitigate the total effect? \textbf{Answer:} NO
\item Does the observed total effect exceed the practical relevance threshold? \textbf{Answer:} YES
\item Does the spurious component (SE) contribute substantially to the observed difference? \textbf{Answer:} NO
\item Is the direct effect the dominant component compared to the indirect effect? \textbf{Answer:} YES
\end{enumerate}

\vspace{1em}
\noindent\footnotesize
The recap answers follow deterministic threshold rules applied to the values above: direct discrimination when $|\mathrm{DE}| > 0.05$; a mediator channel counted as negligible when $|\mathrm{IE}| < 0.005$; practical relevance when $|\mathrm{TV}| > 0.05$; a substantial spurious component when $|\mathrm{SE}| / |\mathrm{TV}| > 0.1$.

\end{document}